\lecture{19}{}{Regular and Singular Matrices}
\section{Regulars}
\begin{note}
    The matrix equation $A\mathbf{x} = \mathbf{b}$ is equivalent to a linear system. A column vector $\mathbf{x}$ is a solution of the matrix equation if and only if the $n$-tuple $x$ is a solution of the linear system.
\end{note}

\subsection{Solving Matrix Equations}

\begin{note}
    Writing $A\mathbf{x} = \mathbf{b}$ instead of the full linear system is not merely convenient notation. It allows us to manipulate equations using matrix algebra. Can we "solve" the system in one step as $\mathbf{x} = \mathbf{b}/A$?
\end{note}

\begin{note}
    Consider the scalar equation $ax = b$ where $a, b \in F$. We can solve for $x$ as $x = b/a$ only if $a \neq 0$. When $a \neq 0$, we have $x = a^{-1}b$. The key property: in a field, every nonzero element has a multiplicative inverse.
\end{note}

\begin{note}
    However, $M_{m \times n}(F)$ is not a field. Even $M_n(F)$ is not a field, though it has better properties:
    \begin{itemize}
        \item Multiplication is always defined and closed: $A, B \in M_n(F) \implies AB \in M_n(F)$
        \item There is a multiplicative neutral element: $I_n \in M_n(F)$ such that $AI_n = I_nA = A$ for all $A \in M_n(F)$
        \item Where there is a neutral element, there may be inverses
    \end{itemize}
\end{note}

\begin{note}
    Not all matrices have inverses:
    \begin{itemize}
        \item The zero matrix $O_n$ has no inverse (similar to how $0$ has no inverse in a field)
        \item Any matrix with a row of zeros has no inverse (if $A$ has a zero row, then $AB$ has a zero row, so $AB \neq I$)
        \item Any matrix with a column of zeros has no inverse (if $A$ has a zero column, then $BA$ has a zero column, so $BA \neq I$)
        \item $\begin{bmatrix} 1 & 1 \\ 1 & 1 \end{bmatrix}$ has no inverse (verify by attempting to find one)
    \end{itemize}
\end{note}

\begin{definition}[Regular and Singular Matrices]\label{def:regular-singular}
    A matrix $A \in M_n(F)$ is called \textbf{regular} (or \textbf{invertible}) if there exists $B \in M_n(F)$ such that
    \[
    AB = BA = I_n
    \]
    In this case, $B$ is called an \textbf{inverse} of $A$. A matrix that is not regular is called \textbf{singular} (or \textbf{non-invertible}).
\end{definition}

\begin{notation}
    When an inverse of $A$ exists, it is unique. We call it \textbf{the inverse} and denote it $A^{-1}$.
\end{notation}

\begin{note}
    Only square matrices can be regular or singular. While you may find $A \in M_{2 \times 3}(F)$ and $B \in M_{3 \times 2}(F)$ such that $AB = I_2$ and $BA = I_3$, these are different identity matrices, so this doesn't constitute invertibility.
\end{note}

\begin{theorem}[One-Sided Inverse Test]\label{thm:one-sided-inverse}
    Let $A, B \in M_n(F)$.
    \begin{enumerate}[label=(\roman*)]
        \item If $A$ is regular and $AB = I$, then $BA = I$
        \item If $A$ is regular and $BA = I$, then $AB = I$
        \item If $AB = I_n$, then $A$ is regular and $B$ is its inverse (proof deferred to Unit 5)
    \end{enumerate}
    In words: if you know that $A$ is regular and suspect that $B$ is its inverse, you only need to verify one of the two products equals $I$.
\end{theorem}

\begin{proof}
    \textbf{(i)}: If $A$ is regular, then $A^{-1}$ exists and
    \[
    BA = I(BA) = (A^{-1}A)(BA) = A^{-1}(AB)A = A^{-1}(I)A = A^{-1}A = I
    \]
    
    \textbf{(ii)}: If $A$ is regular, then $A^{-1}$ exists and
    \[
    AB = (AB)I = (AB)(AA^{-1}) = A(BA)A^{-1} = A(I)A^{-1} = AA^{-1} = I
    \]
\end{proof}

\begin{proposition}[Cancellation Laws]\label{prop:cancellation}
    Let $A \in M_n(F)$ be regular. Then for any $B, C \in M_n(F)$:
    \begin{enumerate}[label=(\roman*)]
        \item $AB = AC \implies B = C$
        \item $BA = CA \implies B = C$
    \end{enumerate}
\end{proposition}

\begin{proof}
    \textbf{(i)}: If $AB = AC$, multiply both sides on the left by $A^{-1}$:
    \[
    A^{-1}(AB) = A^{-1}(AC) \implies (A^{-1}A)B = (A^{-1}A)C \implies IB = IC \implies B = C
    \]
    
    \textbf{(ii)}: If $BA = CA$, multiply both sides on the right by $A^{-1}$:
    \[
    (BA)A^{-1} = (CA)A^{-1} \implies B(AA^{-1}) = C(AA^{-1}) \implies BI = CI \implies B = C
    \]
\end{proof}

\begin{proposition}[Zero Product Properties]\label{prop:zero-product}
    Let $A, B \in M_n(F)$.
    \begin{enumerate}[label=(\roman*)]
        \item If $A$ is regular and $AB = O_n$, then $B = O_n$
        \item If $A, B$ are both regular, then $AB \neq O_n$
        \item If $A \neq O_n$ and $B \neq O_n$ and $AB = O_n$, then $A, B$ are both singular
    \end{enumerate}
\end{proposition}

\begin{proof}
    \textbf{(i)}: If $AB = O_n$, then $AB = AO_n$. By Proposition~\ref{prop:cancellation}(i), $B = O_n$.
    
    \textbf{(ii)}: If both $A$ and $B$ are regular and $AB = O_n$, then by (i), $B = O_n$. But $O_n$ is not regular, contradiction.
    
    \textbf{(iii)}: Suppose $A$ is regular. Then by (i), $B = O_n$, contradicting $B \neq O_n$. Similarly, if $B$ is regular, then $A = O_n$. Therefore both must be singular.
\end{proof}

\begin{theorem}[Properties of Inverses]\label{thm:inverse-properties}
    Let $A, B \in M_n(F)$.
    \begin{enumerate}[label=(\roman*)]
        \item If $A$ is regular, then $A^{-1}$ is also regular and $(A^{-1})^{-1} = A$
        \item If $A$ is regular, then $A^t$ is also regular and $(A^t)^{-1} = (A^{-1})^t$
        \item If $A, B$ are both regular, then $AB$ is also regular and $(AB)^{-1} = B^{-1}A^{-1}$
        \item If $A$ is regular and $s \neq 0 \in F$, then $sA$ is also regular and $(sA)^{-1} = s^{-1}A^{-1}$
    \end{enumerate}
\end{theorem}

\begin{proof}
    \textbf{(i)}: We have $A^{-1}A = AA^{-1} = I$. This shows $A$ is an inverse of $A^{-1}$, so $A^{-1}$ is regular with $(A^{-1})^{-1} = A$.
    
    \textbf{(ii)}: We need to show $(A^{-1})^t A^t = A^t(A^{-1})^t = I$. Using Proposition~\ref{prop:transpose-product}:
    \[
    (A^{-1})^t A^t = (AA^{-1})^t = I^t = I
    \]
    \[
    A^t(A^{-1})^t = (A^{-1}A)^t = I^t = I
    \]
    
    \textbf{(iii)}: We need to show $(AB)(B^{-1}A^{-1}) = (B^{-1}A^{-1})(AB) = I$:
    \[
    (AB)(B^{-1}A^{-1}) = A(BB^{-1})A^{-1} = AIA^{-1} = AA^{-1} = I
    \]
    \[
    (B^{-1}A^{-1})(AB) = B^{-1}(A^{-1}A)B = B^{-1}IB = B^{-1}B = I
    \]
    
    \textbf{(iv)}: We need to show $(sA)(s^{-1}A^{-1}) = (s^{-1}A^{-1})(sA) = I$:
    \[
    (sA)(s^{-1}A^{-1}) = s(A(s^{-1}A^{-1})) = s(s^{-1}(AA^{-1})) = s(s^{-1}I) = (ss^{-1})I = I
    \]
    \[
    (s^{-1}A^{-1})(sA) = s^{-1}(A^{-1}(sA)) = s^{-1}(s(A^{-1}A)) = s^{-1}(sI) = (s^{-1}s)I = I
    \]
\end{proof}

\begin{theorem}[Product of Regular Matrices]\label{thm:product-regular}
    For all $k \in \mathbb{N}$, if $A_1, A_2, \ldots, A_k \in M_n(F)$ are regular, then $A_1A_2\cdots A_k$ is regular and
    \[
    (A_1A_2\cdots A_{k-1}A_k)^{-1} = A_k^{-1}A_{k-1}^{-1}\cdots A_2^{-1}A_1^{-1}
    \]
\end{theorem}

\begin{proof}
    By induction on $k$.
    
    \textbf{Base case ($k=2$)}: This is Theorem~\ref{thm:inverse-properties}(iii).
    
    \textbf{Induction hypothesis}: Suppose for some $k \in \mathbb{N}$, if $A_1, \ldots, A_k \in M_n(F)$ are regular, then $A_1\cdots A_k$ is regular and $(A_1\cdots A_k)^{-1} = A_k^{-1}\cdots A_1^{-1}$.
    
    \textbf{Induction step}: Let $A_1, \ldots, A_k, A_{k+1} \in M_n(F)$ be regular. Let $B = A_1A_2\cdots A_k$. By the induction hypothesis, $B$ is regular and $B^{-1} = A_k^{-1}\cdots A_1^{-1}$. By Theorem~\ref{thm:inverse-properties}(iii), $BA_{k+1}$ is regular and
    \[
    (BA_{k+1})^{-1} = A_{k+1}^{-1}B^{-1} = A_{k+1}^{-1}A_k^{-1}\cdots A_1^{-1}
    \]
    Therefore $A_1\cdots A_kA_{k+1}$ is regular with the claimed inverse.
\end{proof}